\chapter{Myringotomy and Tube Insertion}

\section*{Overview}
Myringotomy with tympanostomy tube placement creates a small opening in the eardrum and inserts a tiny tube to ventilate the middle ear and prevent fluid accumulation.

\section*{Alternative Names}
Ear tube surgery; tympanostomy tube placement; grommet insertion; ventilation tube insertion

\section*{Principle}
Incision of the tympanic membrane allows drainage of middle-ear fluid and ventilation of the middle ear, equalizing pressure. A tube keeps the opening patent, bypassing a dysfunctional eustachian tube.

\section*{History}
Myringotomy has been performed since the 19th century, and tympanostomy tubes were introduced by Armstrong in 1954. They remain the most common pediatric surgical procedure.

\section*{Why It Is Done}
Performed for recurrent acute otitis media, persistent middle-ear effusion (chronic otitis media with effusion) with hearing loss, and eustachian tube dysfunction that has failed medical therapy.

\section*{Is This a Primary or Secondary Test or Procedure}
Primary surgical treatment for recurrent or chronic middle-ear disease; secondary to antibiotic therapy.

\section*{How It Is Performed}
Performed under general anesthesia in children, or local anesthesia in adults. The ear canal is examined microscopically, a small incision is made in the tympanic membrane, fluid is aspirated, and a tube is inserted. The procedure takes only minutes.

\section*{Preparation}
NPO before general anesthesia, and an examination by an ear, nose, and throat specialist. Audiometry is often performed preoperatively.

\section*{Contraindications and Precautions}
Cholesteatoma or a perforation that is already draining is a contraindication. Precaution: the patient should avoid water entering the ear unless protected.

\section*{Risks}
Tube obstruction or premature extrusion, persistent perforation after tube removal, scarring of the eardrum, recurrent drainage (otorrhea), and rarely cholesteatoma.

\section*{Aftercare and Recovery}
Hearing usually improves promptly. The tube stays in place for 6-18 months before extruding spontaneously.

\section*{Outcome and Follow-up}
Resolution of effusion and improved hearing confirm success. Drainage from the ear prompts evaluation for infection.

\section*{Expected Results}
A patent tube with a dry, aerated middle ear and restored hearing.

\section*{Follow-up}
Audiometry and otoscopy at intervals; tubes that fail to extrude or recurrent disease may require repeat placement or removal.

\section*{References}
\begin{enumerate}
  \item MedlinePlus. Myringotomy and tubes. U.S. National Library of Medicine; 2025.
  \item Current Medical Diagnosis and Treatment. McGraw-Hill; 2025.
\end{enumerate}

\clearpage
